\section{Introduction}
\label{sec:intro}

FPGA prototypes allow hardware and system software to evolve before silicon is
available. Their value increasingly depends on running a complete software
stack, including the operating system, production drivers, and realistic I/O
workloads. Such validation requires real NVMe SSDs, NICs, and accelerators:
software device models are configurable and observable, but their usefulness is
bounded by the fidelity of the modeled controller, firmware, timing, and DMA
behavior.

Attaching a real device retains its controller and firmware as participants in
validation, but usually makes the
validation environment rigid. The DUT can use only devices reachable through
the prototype board and its host topology. Adding device types often requires
device-specific RTL for configuration, register access, DMA, and interrupts.
Changing the device composition may require rewiring, additional PCIe hardware,
or a new FPGA build. These costs consume both scarce FPGA resources and
engineering time. We call their incremental cost the \tax.

Moving a device to another host does not by itself remove this coupling. A DUT
does not consume an abstract remote-I/O service. Its unmodified Linux stack
expects a locally discoverable PCIe function with coherent configuration and
BAR state. Its driver creates protocol-specific queues, publishes DMA addresses,
rings doorbells, and assumes that data is visible before a completion or
interrupt arrives. Forwarding MMIO operations without preserving this end-to-end
contract may enumerate a device but cannot make it usable.

This paper asks: \emph{Can the peripheral subsystem visible to an FPGA-hosted
DUT be defined in software, without moving complete device implementations into
the FPGA and without changing the DUT driver?} We answer with \sys, a
software-defined peripheral subsystem for FPGA prototyping. ``Software-defined''
has a precise meaning in \sys: within a hardware-supported capacity, software
selects which logical endpoints the DUT sees, binds them to local or remote
physical backends, and chooses how their control semantics and DMA paths are
realized. The bitstream provides stable mechanisms rather than a fixed device
composition.

The design balances two objectives of peripheral access. Virtualization gives
the user control over the logical device view and backend binding. Passthrough
retains execution by a physical controller, including the firmware and I/O
effects exercised by the workload. These objectives meet at a semantic
boundary: a configurable endpoint still has to issue commands that its backend
can execute, and an adapter can alter the timing and behavior observed by the
driver. \sys uses selective mediation to preserve this boundary, with
functional fidelity and timing fidelity treated as distinct validation goals.

\sys obtains this ability through \emph{dual decoupling}. First,
\emph{device-presentation--backend-implementation decoupling} separates the
PCIe identity and driver-visible behavior presented to the DUT from the
physical device's location and backend realization. Second,
\emph{control-semantics--data-path decoupling} allows software to mediate
configuration and runtime control while bulk data uses a separately optimized
DMA path. These separations create flexibility, but they are safe only if the
resulting distributed state still behaves like one device. \sys therefore
enforces three coherence relations: configuration writes become visible
together with their BAR routes; a submission reaches a backend only after its
queue entry and translated DMA addresses are visible; and a completion becomes
visible only after its data effects.

The FPGA front-end implements standard ECAM access, a configuration shadow,
BAR routing, a control-event transport, DMA windows, and interrupt injection.
The host software back-end owns the logical topology and device state, dispatches
accesses by logical BDF, translates NVMe and NIC queue semantics, and connects
each endpoint to a physical device. This partition keeps latency-critical,
interface-facing mechanisms in hardware while placing extensible device
semantics in software.

This paper makes three contributions:
\begin{itemize}[leftmargin=*]
    \item It defines a software-defined peripheral subsystem that turns the
    DUT-visible device composition and endpoint-to-backend binding into software
    configuration rather than a consequence of physical attachment.
    \item It presents a hardware/software architecture based on device-
    presentation--backend-implementation decoupling and control-semantics--
    data-path decoupling, with explicit coherence rules spanning enumeration,
    runtime control, DMA address translation, completion, and interrupts.
    \item It demonstrates the design on a RISC-V SoC FPGA prototype with real
    NVMe and Intel Ethernet backends, and evaluates native-driver compatibility,
    control overhead, data-path performance, and completion/interrupt ordering.
\end{itemize}

The current implementation targets pre-silicon FPGA prototyping. The same
logical abstraction may be useful in post-silicon validation, but we treat that
extension as future work rather than a demonstrated contribution.
Section~\ref{sec:overview} develops the architectural separation,
Section~\ref{sec:design} derives the mechanisms that preserve device semantics,
and Section~\ref{sec:evaluation} evaluates their compatibility and cost.
Section~\ref{sec:discussion} examines the resulting configurability--fidelity
tradeoff.
